\chapter{Joint Distributions, Independence and Covariance}
\label{ch:joint}

Real problems involve several random quantities at once: the two dice, the
height and weight of a person, the numbers of successes in two groups of trials.
This chapter studies how random variables are related. We define joint and
marginal PMFs, independence of random variables, and the quantities covariance
and correlation that measure linear dependence. We then use them to find
variances of sums and to analyse classical problems.

\section{Joint and marginal distributions}

\begin{definition}[Joint PMF]
For discrete random variables $X,Y$ on the same probability space, the
\emph{joint PMF} is $p_{X,Y}(x,y)=\Prob(X=x,\,Y=y)$.
\end{definition}

The joint PMF is non-negative and sums to $1$ over all pairs $(x,y)$. For any
set $A\subseteq\R^2$, $\Prob((X,Y)\in A)=\sum_{(x,y)\in A}p_{X,Y}(x,y)$.

\begin{proposition}[Marginals]
$p_X(x)=\sum_yp_{X,Y}(x,y)$ and $p_Y(y)=\sum_xp_{X,Y}(x,y)$.
\end{proposition}
\begin{proof}
The events $\{Y=y\}$ partition $\Omega$, so $\{X=x\}=\bigsqcup_y\{X=x,Y=y\}$.
Apply countable additivity.
\end{proof}

\begin{nuance}{marginals do not determine the joint distribution}
Many different joint PMFs have the same marginals. For example, let $X,Y$ each be
Bernoulli$(1/2)$. Taking $Y=X$ gives joint mass $1/2$ on $(0,0)$ and $(1,1)$;
taking $Y=1-X$ gives mass $1/2$ on $(0,1)$ and $(1,0)$; taking independent
$X,Y$ gives mass $1/4$ on each of the four points. The marginals are identical in
all three cases, but the relationships between the variables are completely
different.
\end{nuance}

\begin{theorem}[LOTUS for two variables]\label{thm:lotus2}
If $g:\R^2\to\R$ and $\sum_{x,y}\abs{g(x,y)}p_{X,Y}(x,y)<\infty$, then
$\Exp[g(X,Y)]=\sum_{x,y}g(x,y)p_{X,Y}(x,y)$.
\end{theorem}
\begin{proof}
The proof is the same as for Theorem~\ref{thm:lotus}: group the pairs $(x,y)$
by the value of $g(x,y)$, and regroup the absolutely convergent series.
\end{proof}

\section{Independence of random variables}

\begin{definition}[Independent random variables]
Discrete random variables $X$ and $Y$ are \emph{independent} if the events
$\{X=x\}$ and $\{Y=y\}$ are independent for all $x,y$:
\[
p_{X,Y}(x,y)=p_X(x)\,p_Y(y)\quad\text{for all }x,y.
\]
Random variables $X_1,\dots,X_n$ are \emph{mutually independent} if
$\Prob(X_1=x_1,\dots,X_n=x_n)=\prod_i\Prob(X_i=x_i)$ for all $x_1,\dots,x_n$.
\end{definition}

\begin{proposition}[Independence passes to sets and functions]\label{prop:indfun}
If $X,Y$ are independent discrete random variables then
\begin{enumerate}[label=(\roman*)]
\item $\Prob(X\in A,Y\in B)=\Prob(X\in A)\Prob(Y\in B)$ for all sets $A,B$;
\item $f(X)$ and $g(Y)$ are independent for all functions $f,g$.
\end{enumerate}
\end{proposition}
\begin{proof}
(i) Sum the factorisation over $x\in A$, $y\in B$:
$\sum_{x\in A}\sum_{y\in B}p_X(x)p_Y(y)=\bigl(\sum_{x\in A}p_X(x)\bigr)
\bigl(\sum_{y\in B}p_Y(y)\bigr)$.
(ii) Apply (i) with $A=f^{-1}(\{u\})$ and $B=g^{-1}(\{v\})$:
$\Prob(f(X)=u,g(Y)=v)=\Prob(f(X)=u)\Prob(g(Y)=v)$.
\end{proof}

\begin{theorem}[Expectation of a product]\label{thm:prodexp}
If $X,Y$ are independent and integrable, then $XY$ is integrable and
$\Exp[XY]=\Exp[X]\Exp[Y]$.
\end{theorem}
\begin{proof}
By Theorem~\ref{thm:lotus2} and the factorisation of the joint PMF,
$\Exp[XY]=\sum_{x,y}xy\,p_X(x)p_Y(y)=\bigl(\sum_xxp_X(x)\bigr)\bigl(\sum_yyp_Y(y)\bigr)$,
the rearrangement being justified by absolute convergence
(the same computation with $\abs x\abs y$ shows $\Exp\abs{XY}=\Exp\abs X\Exp\abs Y<\infty$).
\end{proof}

\begin{pitfall}{the converse of Theorem~\ref{thm:prodexp} is false}
$\Exp[XY]=\Exp[X]\Exp[Y]$ does not imply independence. Example~\ref{ex:uncorr}
below shows uncorrelated, strongly dependent variables.
\end{pitfall}

\section{Covariance and correlation}

\begin{definition}[Covariance and correlation]
For random variables with finite second moments, the \emph{covariance} is
\[
\Cov(X,Y)=\Exp\bigl[(X-\mu_X)(Y-\mu_Y)\bigr]=\Exp[XY]-\Exp[X]\Exp[Y],
\]
(the two forms agree by expanding the product and using linearity), and, if both
variances are positive, the \emph{correlation} is
$\rho(X,Y)=\dfrac{\Cov(X,Y)}{\sigma_X\sigma_Y}$. They are \emph{uncorrelated} if
$\Cov(X,Y)=0$.
\end{definition}

\begin{theorem}[Properties of covariance]\label{thm:cov}
\begin{enumerate}[label=(\roman*)]
\item $\Cov(X,X)=\Var(X)$ and $\Cov(X,Y)=\Cov(Y,X)$.
\item $\Cov(aX+b,\,cY+d)=ac\,\Cov(X,Y)$.
\item Bilinearity: $\Cov\bigl(\sum_ia_iX_i,\sum_jb_jY_j\bigr)=\sum_{i,j}a_ib_j\Cov(X_i,Y_j)$.
\item $\Var\bigl(\sum_{i=1}^nX_i\bigr)=\sum_i\Var(X_i)+2\sum_{i<j}\Cov(X_i,X_j)$.
\item If $X,Y$ are independent then $\Cov(X,Y)=0$; hence the variance of a sum of
independent variables is the sum of the variances.
\end{enumerate}
\end{theorem}
\begin{proof}
(i), (ii) are immediate from the definition. (iii) Expand the product inside the
expectation and apply linearity. (iv) By (i) and (iii),
$\Var(\sum X_i)=\Cov(\sum_iX_i,\sum_jX_j)=\sum_{i,j}\Cov(X_i,X_j)$; the diagonal
terms are the variances and the off-diagonal terms appear twice. (v) By
Theorem~\ref{thm:prodexp}.
\end{proof}

\begin{theorem}[Cauchy--Schwarz inequality]\label{thm:cs}
If $\Exp[X^2],\Exp[Y^2]<\infty$ then $\bigl(\Exp[XY]\bigr)^2\le\Exp[X^2]\Exp[Y^2]$,
with equality iff $aX+bY=0$ almost surely for some $(a,b)\ne(0,0)$.
Consequently $\abs{\rho(X,Y)}\le1$, and $\abs\rho=1$ iff $Y$ is an exact linear
function of $X$.
\end{theorem}
\begin{proof}
For every real $t$, $0\le\Exp[(X+tY)^2]=\Exp[X^2]+2t\Exp[XY]+t^2\Exp[Y^2]$. If
$\Exp[Y^2]=0$ then $Y=0$ almost surely and the inequality is trivial; otherwise this
quadratic in $t$ is non-negative for all $t$, so its discriminant is
non-positive: $4(\Exp[XY])^2-4\Exp[X^2]\Exp[Y^2]\le0$. Equality means the
quadratic has a double root $t_0$, i.e.\ $\Exp[(X+t_0Y)^2]=0$ and so
$X+t_0Y=0$ a.s. Applying the inequality to $X-\mu_X$ and $Y-\mu_Y$ gives
$\Cov(X,Y)^2\le\Var(X)\Var(Y)$, i.e.\ $\abs\rho\le1$.
\end{proof}

\begin{example}[Uncorrelated but dependent]\label{ex:uncorr}
Let $X$ be uniform on $\{-1,0,1\}$ and $Y=X^2$. Then $\Exp[X]=0$ and
$\Exp[XY]=\Exp[X^3]=0$, so $\Cov(X,Y)=0$. But $Y$ is a function of $X$: for
example $\Prob(X=0,Y=0)=1/3\ne\Prob(X=0)\Prob(Y=0)=1/9$.
\end{example}

\begin{nuance}{correlation measures \emph{linear} association only}
Correlation zero says that no linear relationship exists; it does not say the
variables have no relationship at all. In Example~\ref{ex:uncorr} the dependence
is perfect but non-monotone. For jointly Gaussian variables (Chapter~\ref{ch:jointcont})
uncorrelated does imply independent; this is a special feature of the Gaussian
family.
\end{nuance}

\section{Sums of independent random variables}

\begin{theorem}[Convolution]\label{thm:convolution}
If $X,Y$ are independent integer-valued random variables, then $Z=X+Y$ has PMF
\[
p_Z(z)=\sum_xp_X(x)\,p_Y(z-x).
\]
\end{theorem}
\begin{proof}
The event $\{Z=z\}$ is the disjoint union over $x$ of
$\{X=x,Y=z-x\}$; apply independence to each term.
\end{proof}

\begin{theorem}[Closure under sums]\label{thm:closure}
Let $X,Y$ be independent.
\begin{enumerate}[label=(\roman*)]
\item If $X\sim\mathrm{Bin}(m,p)$ and $Y\sim\mathrm{Bin}(n,p)$ (same $p$), then
$X+Y\sim\mathrm{Bin}(m+n,p)$.
\item If $X\sim\mathrm{Poisson}(\lambda)$ and $Y\sim\mathrm{Poisson}(\mu)$, then
$X+Y\sim\mathrm{Poisson}(\lambda+\mu)$.
\end{enumerate}
\end{theorem}
\begin{proof}
(i) By convolution and Vandermonde's identity (Theorem~\ref{thm:vander}),
\[
\Prob(X+Y=z)=\sum_k\binom mkp^kq^{m-k}\binom n{z-k}p^{z-k}q^{n-z+k}
=p^zq^{m+n-z}\sum_k\binom mk\binom n{z-k}=\binom{m+n}zp^zq^{m+n-z}.
\]
(ii) By convolution and the binomial theorem,
\[
\Prob(X+Y=z)=\sum_{k=0}^ze^{-\lambda}\frac{\lambda^k}{k!}e^{-\mu}\frac{\mu^{z-k}}{(z-k)!}
=\frac{e^{-(\lambda+\mu)}}{z!}\sum_{k=0}^z\binom zk\lambda^k\mu^{z-k}
=e^{-(\lambda+\mu)}\frac{(\lambda+\mu)^z}{z!}.\qedhere
\]
\end{proof}

Part (i) matches the story: $X+Y$ counts successes in $m+n$ independent trials.
Part (ii) matches the Poisson limit: merging two independent streams of rare events
gives a stream with the total rate.

\begin{example}[Variance of the binomial via independence]
$X\sim\mathrm{Bin}(n,p)$ is a sum of $n$ independent Bernoulli$(p)$ variables, each
with variance $pq$. By Theorem~\ref{thm:cov}(v), $\Var(X)=npq$, matching
Theorem~\ref{thm:binmoments}. Similarly $\mathrm{NB}(r,p)$ is a sum of $r$
independent Geom$(p)$ variables (the waits between successes), so its variance
is $r\cdot q/p^2$.
\end{example}

\section{Worked examples}

\begin{example}[Variance of the hypergeometric]\label{ex:hypervar}
Let $X=\sum_{j=1}^nI_j$ as before, with $p=K/N$ and $q=1-p$. Then
$\Var(I_j)=pq$. For $i\ne j$,
\[
\Prob(I_i=I_j=1)=\frac KN\cdot\frac{K-1}{N-1},
\]
so
\[
\Cov(I_i,I_j)=\frac{K(K-1)}{N(N-1)}-\Bigl(\frac KN\Bigr)^2
=p\Bigl[\frac{K-1}{N-1}-\frac KN\Bigr]=p\,\frac{N(K-1)-K(N-1)}{N(N-1)}
=-\frac{pq}{N-1}.
\]
Using Theorem~\ref{thm:cov}(iv):
\[
\Var(X)=npq+n(n-1)\Bigl(-\frac{pq}{N-1}\Bigr)=npq\Bigl[1-\frac{n-1}{N-1}\Bigr]
=npq\,\frac{N-n}{N-1}.
\]
The covariance is negative: drawing a red ball makes red slightly less likely
next time. When $N\to\infty$ the covariance vanishes and we recover the binomial.
\end{example}

\begin{example}[Coupon collector]\label{ex:coupon}
There are $n$ types of coupon, each new coupon uniformly random and independent.
Let $T$ be the number of coupons collected until all $n$ types have appeared. Write
$T=T_1+\dots+T_n$ where $T_i$ is the number of additional coupons needed to
get the $i$th new type after having $i-1$ types. While we hold $i-1$ types, each
draw gives a new type with probability $p_i=(n-i+1)/n$, so $T_i\sim\mathrm{Geom}(p_i)$.
By linearity,
\[
\Exp[T]=\sum_{i=1}^n\frac1{p_i}=\sum_{i=1}^n\frac n{n-i+1}=nH_n\approx n(\ln n+\gamma).
\]
The $T_i$ are independent (the number of draws needed at each stage depends only
on fresh draws), so
$\Var(T)=\sum_i\frac{1-p_i}{p_i^2}=\sum_i\frac{n(i-1)}{(n-i+1)^2}\le\frac{\pi^2}6n^2$.
For $n=52$ cards (drawing with replacement) the mean is
$52H_{52}\approx236$ draws.
\end{example}

\begin{example}[Multinomial distribution]\label{ex:multinomial}
Perform $n$ independent trials, each resulting in category $i\in\{1,\dots,r\}$ with
probability $p_i$ ($\sum p_i=1$). Let $X_i$ be the number of trials in category
$i$. Counting sequences (Theorem~\ref{thm:multinomial}),
\[
\Prob(X_1=n_1,\dots,X_r=n_r)=\frac{n!}{n_1!\cdots n_r!}p_1^{n_1}\cdots p_r^{n_r},
\qquad n_1+\dots+n_r=n .
\]
Each marginal $X_i$ is $\mathrm{Bin}(n,p_i)$. For $i\ne j$ write
$X_i=\sum_t\indic_{\{\text{trial }t\text{ in }i\}}$ and
$X_j=\sum_s\indic_{\{\text{trial }s\text{ in }j\}}$. Indicators from different
trials are independent, so only $s=t$ contribute to the covariance; in the same
trial the indicators cannot both be $1$:
$\Cov=\sum_t(0-p_ip_j)=-np_ip_j$. The counts are negatively correlated: a larger
count in one category leaves fewer trials for the others.
\end{example}

\begin{interview}{covariance of indicators}
Variances of counts that are sums of dependent indicators are computed with
$\Var(\sum I_j)=\sum\Var(I_j)+\sum_{i\ne j}\Cov(I_i,I_j)$, and
$\Cov(I_i,I_j)=\Prob(\text{both})-\Prob(I_i)\Prob(I_j)$. The only new work is the
joint probability of pairs. Typical applications: fixed points in a permutation
(variance is $1$ for $n\ge2$; check this), records, runs, edges in a random graph.
\end{interview}

\section*{Exercises}
\addcontentsline{toc}{section}{Exercises}

\begin{enumerate}
\item Two fair dice give $X$ (first) and $M=\max$. Find the joint PMF table, the
marginals, and $\Cov(X,M)$.
\item Prove that if $X,Y$ are independent then
$\Var(X-Y)=\Var(X)+\Var(Y)$.
\item Prove $\Cov(X+Y,X-Y)=\Var(X)-\Var(Y)$.
\item Show that for a uniformly random permutation of $\{1,\dots,n\}$ the number
of fixed points has variance $1$ for $n\ge2$. (Indicators; $\Prob(\text{both
fixed})=1/(n(n-1))$.)
\item Let $X,Y$ be independent Geom$(p)$. Find the PMF of $X+Y$ and recognise it
as a negative binomial.
\item Prove that if $X,Y$ are independent Poisson variables then, conditionally on
$X+Y=n$, $X\sim\mathrm{Bin}(n,\lambda/(\lambda+\mu))$.
\item (Collisions) $m$ balls are thrown into $n$ boxes uniformly and
independently. Find the expected number of empty boxes and of pairs of balls in
the same box. For which $m$ (in terms of $n$) is the expected number of colliding
pairs about $1$?
\item In the coupon collector problem find the exact value of $\Var(T)$ for $n=2$
and verify it using the formula.
\item (Challenge) Prove that $\abs{\rho(X,Y)}=1$ implies $Y=aX+b$ almost surely,
and that the sign of $a$ equals the sign of $\rho$.
\end{enumerate}
